\documentclass[10pt]{beamer}

\usetheme{Madrid}

\usepackage{graphicx}

\usepackage{booktabs}

\title{Uncertainty from $b$ quark masses in Z pT spectrum}

\subtitle{Status summary}

\author{AlphaS from Z pT analysis team, CMS}

\date{February 18, 2026}

\begin{document}

\frame{\titlepage}

\begin{frame}{Overview: Samples, Goal, Methodology}
\begin{itemize}
  \item Goal: build a $b$-mass uncertainty to apply to our simulation, since our nominal MiNNLO treats $b$'s as massless.
  \begin{itemize}
    \item N.B.: we are already accounting for heavy quark mass effects in PDFs.
  \end{itemize}
  \item Method:
  \begin{itemize}
    \item select events affected by heavy $b$ quark mass.
    \item Build a corrected distribution as massless-$b$ MiNNLO (unselected) + massive-$b$ MiNNLO (selected).
    \item The difference between the corrected distribution and the nominal MiNNLO is taken as an uncertainty.
  \end{itemize}
  \item We attempted to use the B hadrons to define the set of events that should be swapped, but it doesn't seem to work nicely. We looked at jets and LHE quarks as well, see backup.
  \item Samples:
  \begin{itemize}
    \item nominal MiNNLO: massless $b$, 5FS
    \item dedicated heavy $b$ quark Z$b\bar{b}$ MiNNLO: massive $b$, 4FS. (From arXiv2404.08598: $\sigma=25.28$pb from Table 1, but unclear if that is inclusive $\ell=e,\mu$ or for each lepton flavor.)
  \end{itemize}
\end{itemize}
\end{frame}

\begin{frame}{B-Hadron Diagnostics}
\begin{itemize}
  \item B hadrons are identified from GenPart using PDG-ID hadron-content logic (contains bottom flavor) with GenPart status 1 or 2.
  \begin{itemize}
    \item Exact implementation is shown in backup.
  \end{itemize}
\end{itemize}
\begin{columns}[T,onlytextwidth]
\column{0.5\textwidth}
\includegraphics[width=0.95\linewidth]{\detokenize{/home/submit/lavezzo/public_html/alphaS/260218_z_bb/hadrons/bhad_hiStat_260218_200thr_genmother_iter2_m91w30_refresh/samples_comparison_leadB_pt.png}}
\column{0.5\textwidth}
\includegraphics[width=0.95\linewidth]{\detokenize{/home/submit/lavezzo/public_html/alphaS/260218_z_bb/hadrons/bhad_hiStat_260218_200thr_genmother_iter2_m91w30_refresh/samples_comparison_nBhad.png}}
\end{columns}
\begin{itemize}
  \item \small Large, unexplained normalization difference, on top of what may be a reasonable shape difference at low $p_T$.
\end{itemize}
\end{frame}

\begin{frame}{Swap Result and Definition}
\begin{itemize}
  \item Swap definition: select events with at least one such $B$ hadron (no $p_T$ selection); corrected distribution is built as nominal MiNNLO (unselected) + massive-$b$ MiNNLO (selected).
\end{itemize}
\begin{center}
\includegraphics[width=0.62\textwidth,height=0.42\textheight,keepaspectratio]{\detokenize{/home/submit/lavezzo/public_html/alphaS/260218_z_bb/hadrons/bhad_hiStat_260218_200thr_genmother_iter2_m91w30_refresh/inclusive_ptVgen.png}}
\end{center}
\begin{itemize}
  \item \small Result: swapping gives a sizable normalization shift (not shape-only), so it is too aggressive as-is.
  \item \small Questions: is this due to correct normalization difference between samples? Are we selecting more events from the nominal MiNNLO than the ones we are putting back in from the massive-$b$ MiNNLO? Or is it something else?
\end{itemize}
\end{frame}

\begin{frame}{Backup}

\end{frame}

\begin{frame}{Backup: LHE Multiplicity}
\begin{itemize}
  \item \small LHE $b$-quark observables built with $|\mathrm{pdgId}|=5$ and final-state status.
\end{itemize}
\begin{columns}[T,onlytextwidth]
\column{0.5\textwidth}
\includegraphics[width=0.95\linewidth]{\detokenize{/home/submit/lavezzo/public_html/alphaS/260218_z_bb/hadrons/bhad_hiStat_260218_200thr_genmother_iter2_m91w30_refresh/samples_comparison_n_lhe_init_bbbar.png}}
\column{0.5\textwidth}
\includegraphics[width=0.95\linewidth]{\detokenize{/home/submit/lavezzo/public_html/alphaS/260218_z_bb/hadrons/bhad_hiStat_260218_200thr_genmother_iter2_m91w30_refresh/samples_comparison_n_lhe_fin_bbbar.png}}
\end{columns}
\end{frame}

\begin{frame}{Backup: LHE Leading-$b$ $p_T$ and Total Multiplicity}
\begin{itemize}
  \item \small LHE $b$-quark observables built with $|\mathrm{pdgId}|=5$ and final-state status.
\end{itemize}
\begin{columns}[T,onlytextwidth]
\column{0.5\textwidth}
\includegraphics[width=0.95\linewidth]{\detokenize{/home/submit/lavezzo/public_html/alphaS/260218_z_bb/hadrons/bhad_hiStat_260218_200thr_genmother_iter2_m91w30_refresh/samples_comparison_lhe_bbbar_fin_max_pt.png}}
\column{0.5\textwidth}
\includegraphics[width=0.95\linewidth]{\detokenize{/home/submit/lavezzo/public_html/alphaS/260218_z_bb/hadrons/bhad_hiStat_260218_200thr_genmother_iter2_m91w30_refresh/samples_comparison_n_lhe_bbbar.png}}
\end{columns}
\end{frame}

\begin{frame}{Backup: LHE $b\bar{b}$ Kinematics}
\begin{itemize}
  \item \small LHE $b$-quark observables built with $|\mathrm{pdgId}|=5$ and final-state status.
\end{itemize}
\begin{columns}[T,onlytextwidth]
\column{0.5\textwidth}
\includegraphics[width=0.95\linewidth]{\detokenize{/home/submit/lavezzo/public_html/alphaS/260216_z_bb/hadrons/bhad_dualdef_260216_50thr_quick_rrfix/samples_comparison_m_bb_lhe.png}}
\column{0.5\textwidth}
\includegraphics[width=0.95\linewidth]{\detokenize{/home/submit/lavezzo/public_html/alphaS/260216_z_bb/hadrons/bhad_dualdef_260216_50thr_quick_rrfix/samples_comparison_dR_bb_lhe.png}}
\end{columns}
\end{frame}

\begin{frame}{Backup: GenJet $b$-jet Observables}
\begin{itemize}
  \item \small Gen $b$-jet observables here use hadronFlavour = 5 with no kinematic requirement ($p_T > 0$, $|\eta| < 10$).
\end{itemize}
\begin{columns}[T,onlytextwidth]
\column{0.5\textwidth}
\includegraphics[width=0.95\linewidth]{\detokenize{/home/submit/lavezzo/public_html/alphaS/260218_z_bb/hadrons/bhad_hiStat_260218_200thr_genmother_iter2_m91w30_refresh/samples_comparison_n_bjets.png}}
\column{0.5\textwidth}
\includegraphics[width=0.95\linewidth]{\detokenize{/home/submit/lavezzo/public_html/alphaS/260218_z_bb/hadrons/bhad_hiStat_260218_200thr_genmother_iter2_m91w30_refresh/samples_comparison_lead_bjet_pt.png}}
\end{columns}
\end{frame}

\begin{frame}{Backup: GenJet $\Delta R_{bb}$ and $m_{bb}$}
\begin{itemize}
  \item \small Jet-level pair observables from gen $b$-jets, with the same no-cut jet definition (hadronFlavour = 5, $p_T > 0$, $|\eta| < 10$).
\end{itemize}
\begin{columns}[T,onlytextwidth]
\column{0.5\textwidth}
\includegraphics[width=0.95\linewidth]{\detokenize{/home/submit/lavezzo/public_html/alphaS/260218_z_bb/hadrons/bhad_hiStat_260218_200thr_genmother_iter2_m91w30_refresh/samples_comparison_dR_bb_jet.png}}
\column{0.5\textwidth}
\includegraphics[width=0.95\linewidth]{\detokenize{/home/submit/lavezzo/public_html/alphaS/260218_z_bb/hadrons/bhad_hiStat_260218_200thr_genmother_iter2_m91w30_refresh/samples_comparison_m_bb_jet.png}}
\end{columns}
\end{frame}

\begin{frame}[fragile]{Backup: Exact B-hadron Selection Code}
\begin{block}{Implementation}
\begin{verbatim}
// wrem::finalStateBHadronIdx implementation
if (status[i] != 1 && status[i] != 2) continue;
const int apdg = std::abs(pdgId[i]);
if (!ThePEG::PDT::hasBottom(apdg)) continue;
idx.push_back((int)i);

// Histmaker usage
df = df.Define("bHadIdx", "wrem::finalStateBHadronIdx(GenPart_pdgId, GenPart_status)")
df = df.Define("bHad_pt", "Take(GenPart_pt, bHadIdx)")
df = df.Define("nBhad", "static_cast<int>(bHad_pt.size())")
df = df.Define("bottom_sel", "(bHad_pt.size() >= 1)")
\end{verbatim}
\end{block}
\end{frame}

\begin{frame}{Backup: Unnormalized Swap with 4FS Scaled by 2}
\begin{itemize}
  \item \small Same unnormalized swap definition as main slide, but with an extra factor-2 scaling applied to the selected 4FS component before replacement.
\end{itemize}
\begin{center}
\includegraphics[width=0.62\textwidth,height=0.42\textheight,keepaspectratio]{\detokenize{/home/submit/lavezzo/public_html/alphaS/260218_z_bb/hadrons/bhad_hiStat_260218_200thr_genmother_iter2_m91w30_x2/inclusive_ptVgen_x2.png}}
\end{center}
\end{frame}

\end{document}
